\subsection{Ordered Component Storage and Positional Access}

A sequence is an object that stores components in a fixed left-to-right order. The word ``fixed'' does not mean that the object itself is always immutable. A list can still be changed. It means that, at any given moment, the components have positions.

Python has several built-in sequence types. In this chapter, the main ones are:

\begin{itemize}
    \item \texttt{str}: a sequence of text characters;
    \item \texttt{list}: a mutable sequence of object references;
    \item \texttt{tuple}: an immutable sequence of object references.
\end{itemize}

The same basic access rules apply to all of them: indexing, slicing, length testing, membership testing, and iteration.

\begin{verbatim}
text = "D'oh!"
items = ["spam", "eggs", 42, "larch"]
point = (10, 20, 30)

print(f"text:  {text!r}")
print(f"items: {items!r}")
print(f"point: {point!r}")

print()
print(f"type(text):  {type(text)}")
print(f"type(items): {type(items)}")
print(f"type(point): {type(point)}")
\end{verbatim}

Possible output:

\begin{verbatim}
text:  "D'oh!"
items: ['spam', 'eggs', 42, 'larch']
point: (10, 20, 30)

type(text):  <class 'str'>
type(items): <class 'list'>
type(point): <class 'tuple'>
\end{verbatim}

At the object level, these are different types. At the sequence level, they share a common positional model.

A first inspection with \texttt{dir()} shows that these objects expose many operations. The full output of \texttt{dir()} is long, so the following program checks only a few names that matter for basic sequence behavior.

\begin{verbatim}
text = "D'oh!"
items = ["spam", "eggs", 42, "larch"]
point = (10, 20, 30)

names = ["__len__", "__getitem__", "__contains__", "index", "count"]

print("selected sequence-related attributes")
print()

for obj_name, obj in [("text", text), ("items", items), ("point", point)]:
    print(f"{obj_name}: {type(obj)}")
    for name in names:
        print(f"  {name:<14} {name in dir(obj)}")
    print()
\end{verbatim}

Possible output:

\begin{verbatim}
selected sequence-related attributes

text: <class 'str'>
  __len__        True
  __getitem__    True
  __contains__   True
  index          True
  count          True

items: <class 'list'>
  __len__        True
  __getitem__    True
  __contains__   True
  index          True
  count          True

point: <class 'tuple'>
  __len__        True
  __getitem__    True
  __contains__   True
  index          True
  count          True
\end{verbatim}

The names beginning and ending with double underscores are special method names. They are not usually called directly in beginner-level code. Instead, Python calls them when normal syntax is used:

\begin{itemize}
    \item \texttt{len(obj)} uses the object's length behavior;
    \item \texttt{obj[i]} uses the object's item-access behavior;
    \item \texttt{x in obj} uses the object's membership behavior.
\end{itemize}

So the syntax is simple, but the action still belongs to the object.

\subsubsection{Indexing: Zero-Based Component Addressing with Positive and Negative Offsets}

Indexing extracts one component from a sequence. Python indexes start at zero, not one. The first component is at position \texttt{0}, the second at position \texttt{1}, and so on.

\begin{verbatim}
quote = "Ni!"

print(f"quote: {quote!r}")
print(f"quote[0]: {quote[0]!r}")
print(f"quote[1]: {quote[1]!r}")
print(f"quote[2]: {quote[2]!r}")
\end{verbatim}

Possible output:

\begin{verbatim}
quote: 'Ni!'
quote[0]: 'N'
quote[1]: 'i'
quote[2]: '!'
\end{verbatim}

For a string, indexing returns a one-character string object. Python does not have a separate character type comparable to C's \texttt{char}.

\begin{verbatim}
quote = "Ni!"
ch = quote[0]

print(f"ch: {ch!r}")
print(f"type(ch): {type(ch)}")
\end{verbatim}

Possible output:

\begin{verbatim}
ch: 'N'
type(ch): <class 'str'>
\end{verbatim}

The same positional rule applies to lists and tuples.

\begin{verbatim}
items = ["Kramer", "Jerry", "Elaine", "George"]
point = (42, 100, 200)

print(f"items[0]: {items[0]}")
print(f"items[1]: {items[1]}")
print(f"items[2]: {items[2]}")
print(f"items[3]: {items[3]}")

print()
print(f"point[0]: {point[0]}")
print(f"point[1]: {point[1]}")
print(f"point[2]: {point[2]}")
\end{verbatim}

Possible output:

\begin{verbatim}
items[0]: Kramer
items[1]: Jerry
items[2]: Elaine
items[3]: George

point[0]: 42
point[1]: 100
point[2]: 200
\end{verbatim}

A positive index counts from the left. A negative index counts from the right.

\begin{verbatim}
line = "Nobody expects the Spanish Inquisition"

print(f"line[0]:  {line[0]!r}")
print(f"line[1]:  {line[1]!r}")
print(f"line[-1]: {line[-1]!r}")
print(f"line[-2]: {line[-2]!r}")
\end{verbatim}

Possible output:

\begin{verbatim}
line[0]:  'N'
line[1]:  'o'
line[-1]: 'n'
line[-2]: 'o'
\end{verbatim}

Negative indexing is not a different storage direction. It is a convenience rule. Python translates the negative offset into a real position from the left.

For example, in a four-element list:

\begin{verbatim}
names = ["Homer", "Marge", "Bart", "Lisa"]

print(f"names[0]:  {names[0]}")
print(f"names[1]:  {names[1]}")
print(f"names[2]:  {names[2]}")
print(f"names[3]:  {names[3]}")

print()
print(f"names[-1]: {names[-1]}")
print(f"names[-2]: {names[-2]}")
print(f"names[-3]: {names[-3]}")
print(f"names[-4]: {names[-4]}")
\end{verbatim}

Possible output:

\begin{verbatim}
names[0]:  Homer
names[1]:  Marge
names[2]:  Bart
names[3]:  Lisa

names[-1]: Lisa
names[-2]: Bart
names[-3]: Marge
names[-4]: Homer
\end{verbatim}

The valid index range depends on the sequence length. For a sequence of length \texttt{4}, the valid positive indexes are \texttt{0}, \texttt{1}, \texttt{2}, and \texttt{3}. The valid negative indexes are \texttt{-1}, \texttt{-2}, \texttt{-3}, and \texttt{-4}.

An index outside the valid range raises an error.

\begin{verbatim}
names = ["Homer", "Marge", "Bart", "Lisa"]

print(names[4])
\end{verbatim}

Possible output:

\begin{verbatim}
IndexError: list index out of range
\end{verbatim}

The error is useful. Python does not return random memory content. It refuses the invalid access.

\subsubsection{Slicing: Extracting Sub-Sequences with \texttt{[start:stop:step]}}

Slicing extracts a new sequence from part of an existing sequence. The general form is:

\begin{verbatim}
sequence[start:stop:step]
\end{verbatim}

The important rule is that \texttt{start} is included, but \texttt{stop} is excluded.

\begin{verbatim}
text = "abcdef"

print(f"text:       {text!r}")
print(f"text[0:3]:  {text[0:3]!r}")
print(f"text[2:5]:  {text[2:5]!r}")
\end{verbatim}

Possible output:

\begin{verbatim}
text:       'abcdef'
text[0:3]:  'abc'
text[2:5]:  'cde'
\end{verbatim}

The slice \texttt{text[0:3]} means: start at index \texttt{0}, move up to index \texttt{3}, but do not include index \texttt{3}.

\begin{verbatim}
text = "abcdef"

print(f"index 0: {text[0]!r}")
print(f"index 1: {text[1]!r}")
print(f"index 2: {text[2]!r}")
print(f"index 3: {text[3]!r}")

print()
print(f"text[0:3]: {text[0:3]!r}")
\end{verbatim}

Possible output:

\begin{verbatim}
index 0: 'a'
index 1: 'b'
index 2: 'c'
index 3: 'd'

text[0:3]: 'abc'
\end{verbatim}

If \texttt{start} is omitted, Python starts at the beginning. If \texttt{stop} is omitted, Python continues to the end.

\begin{verbatim}
quote = "These pretzels are making me thirsty"

print(f"quote[:5]:   {quote[:5]!r}")
print(f"quote[6:14]: {quote[6:14]!r}")
print(f"quote[15:]:  {quote[15:]!r}")
\end{verbatim}

Possible output:

\begin{verbatim}
quote[:5]:   'These'
quote[6:14]: 'pretzels'
quote[15:]:  'are making me thirsty'
\end{verbatim}

The same slicing model applies to lists.

\begin{verbatim}
items = ["spam", "eggs", "bacon", "beans", 42]

print(f"items:      {items}")
print(f"items[:2]:  {items[:2]}")
print(f"items[2:4]: {items[2:4]}")
print(f"items[3:]:  {items[3:]}")
\end{verbatim}

Possible output:

\begin{verbatim}
items:      ['spam', 'eggs', 'bacon', 'beans', 42]
items[:2]:  ['spam', 'eggs']
items[2:4]: ['bacon', 'beans']
items[3:]:  ['beans', 42]
\end{verbatim}

The result of slicing has the same sequence type as the original sequence.

\begin{verbatim}
text = "abcdef"
items = ["a", "b", "c", "d"]
point = (10, 20, 30, 40)

text_part = text[1:3]
items_part = items[1:3]
point_part = point[1:3]

print(f"text_part:  {text_part!r:<20} type: {type(text_part)}")
print(f"items_part: {items_part!r:<20} type: {type(items_part)}")
print(f"point_part: {point_part!r:<20} type: {type(point_part)}")
\end{verbatim}

Possible output:

\begin{verbatim}
text_part:  'bc'                 type: <class 'str'>
items_part: ['b', 'c']            type: <class 'list'>
point_part: (20, 30)              type: <class 'tuple'>
\end{verbatim}

The optional \texttt{step} controls how many positions Python jumps after taking a component.

\begin{verbatim}
text = "abcdefghij"

print(f"text[0:10:1]: {text[0:10:1]!r}")
print(f"text[0:10:2]: {text[0:10:2]!r}")
print(f"text[1:10:2]: {text[1:10:2]!r}")
\end{verbatim}

Possible output:

\begin{verbatim}
text[0:10:1]: 'abcdefghij'
text[0:10:2]: 'acegi'
text[1:10:2]: 'bdfhj'
\end{verbatim}

A negative step walks backward.

\begin{verbatim}
text = "abcdef"

print(f"text[::-1]:  {text[::-1]!r}")
print(f"text[4:1:-1]: {text[4:1:-1]!r}")
\end{verbatim}

Possible output:

\begin{verbatim}
text[::-1]:  'fedcba'
text[4:1:-1]: 'edc'
\end{verbatim}

The compact form \texttt{text[::-1]} is often used to create a reversed copy of a sequence.

Slicing is safer than indexing when boundaries are outside the sequence. Invalid indexing raises an error, but slicing trims the requested range to the available components.

\begin{verbatim}
text = "abc"

print(f"text[0:50]: {text[0:50]!r}")
print(f"text[20:30]: {text[20:30]!r}")
\end{verbatim}

Possible output:

\begin{verbatim}
text[0:50]: 'abc'
text[20:30]: ''
\end{verbatim}

This does not mean slicing ignores all mistakes. A step of zero is impossible because Python would not know how to move through the sequence.

\begin{verbatim}
text = "abc"

print(text[0:3:0])
\end{verbatim}

Possible output:

\begin{verbatim}
ValueError: slice step cannot be zero
\end{verbatim}

\subsubsection{Length, Membership, and Iteration: \texttt{len()}, \texttt{in}, and Sequential Traversal}

The function \texttt{len()} returns the number of components in a sequence.

\begin{verbatim}
text = "D'oh!"
items = ["spam", "eggs", 42]
point = (10, 20, 30, 40)

print(f"len(text):  {len(text)}")
print(f"len(items): {len(items)}")
print(f"len(point): {len(point)}")
\end{verbatim}

Possible output:

\begin{verbatim}
len(text):  5
len(items): 3
len(point): 4
\end{verbatim}

For a string, \texttt{len()} counts text components, not printed screen width and not storage bytes. The byte representation issue belongs to the later text-encoding section.

The operator \texttt{in} tests membership.

\begin{verbatim}
text = "Nobody expects the Spanish Inquisition"
items = ["spam", "eggs", 42, "larch"]
point = (10, 20, 30)

print(f"{'Spanish' in text = }")
print(f"{'French' in text = }")

print()
print(f"{42 in items = }")
print(f"{'eggs' in items = }")
print(f"{'beans' in items = }")

print()
print(f"{20 in point = }")
print(f"{99 in point = }")
\end{verbatim}

Possible output:

\begin{verbatim}
'Spanish' in text = True
'French' in text = False

42 in items = True
'eggs' in items = True
'beans' in items = False

20 in point = True
99 in point = False
\end{verbatim}

For strings, \texttt{in} can test whether a smaller string occurs inside a larger string. For lists and tuples, \texttt{in} tests whether an element equal to the searched value exists as one component.

\begin{verbatim}
text = "abc"
items = ["abc", "def"]

print(f"{'b' in text = }")
print(f"{'b' in items = }")
print(f"{'abc' in items = }")
\end{verbatim}

Possible output:

\begin{verbatim}
'b' in text = True
'b' in items = False
'abc' in items = True
\end{verbatim}

This distinction matters. The string \texttt{"abc"} contains the smaller string \texttt{"b"}. The list \texttt{["abc", "def"]} does not contain \texttt{"b"} as a separate list element.

Iteration means visiting the components of a sequence one at a time.

\begin{verbatim}
names = ["Homer", "Marge", "Bart", "Lisa"]

for name in names:
    print(f"current name: {name}")
\end{verbatim}

Possible output:

\begin{verbatim}
current name: Homer
current name: Marge
current name: Bart
current name: Lisa
\end{verbatim}

The same idea works for strings. Iterating over a string visits one-character string objects.

\begin{verbatim}
word = "Ni!"

for ch in word:
    print(f"ch={ch!r:<4} type={type(ch)}")
\end{verbatim}

Possible output:

\begin{verbatim}
ch='N'  type=<class 'str'>
ch='i'  type=<class 'str'>
ch='!'  type=<class 'str'>
\end{verbatim}

When both the position and the component are needed, \texttt{enumerate()} is clearer than manually maintaining a counter.

\begin{verbatim}
names = ["Jerry", "Elaine", "George", "Kramer"]

for index, name in enumerate(names):
    print(f"{index:>2}: {name}")
\end{verbatim}

Possible output:

\begin{verbatim}
 0: Jerry
 1: Elaine
 2: George
 3: Kramer
\end{verbatim}

A loop can also accumulate a result. The following program walks through a list of numbers and computes their sum. This is intentionally simple: the point is traversal, not a complicated mathematical formula.

\begin{verbatim}
numbers = [10, 20, 12]
total = 0

for nr in numbers:
    print(f"before: total={total:<3} nr={nr}")
    total = total + nr
    print(f"after:  total={total:<3}")
    print()

print(f"final total: {total}")
\end{verbatim}

Possible output:

\begin{verbatim}
before: total=0   nr=10
after:  total=10

before: total=10  nr=20
after:  total=30

before: total=30  nr=12
after:  total=42

final total: 42
\end{verbatim}

The loop does not need to know whether the sequence is backed by a string object, a list object, or a tuple object. It asks the object for its next component according to Python's iteration protocol.

The practical rule for this section is:

\begin{quote}
A sequence is an ordered object. Indexing selects one component. Slicing creates a smaller sequence. \texttt{len()} counts components. \texttt{in} tests containment. A \texttt{for} loop traverses the components from left to right.
\end{quote}